\documentclass[12pt,a4paper]{article}

% ---- Encoding & Font ----
\usepackage{iftex}
\ifPDFTeX
  \usepackage[utf8]{inputenc}
  \usepackage[T1]{fontenc}
\else
  \usepackage{fontspec}  % XeTeX/LuaTeX (tectonic): Unicode nativo (·, —, ¿, ª, §)
\fi
\usepackage[spanish]{babel}
\usepackage{csquotes}

% ---- Page layout ----
\usepackage[top=2.5cm, bottom=2.5cm, left=2.5cm, right=2.5cm]{geometry}
\usepackage{setspace}
\onehalfspacing

% ---- Math ----
\usepackage{amsmath, amssymb, amsthm}
\usepackage{mathtools}
\usepackage{physics}
\usepackage{esint}  % \oiint

% ---- Graphics ----
\usepackage{graphicx}
\usepackage{tikz}
\usepackage{float}
\usepackage{caption}

% ---- Tables ----
\usepackage{array}
\usepackage{booktabs}
\usepackage{tabularx}
\usepackage{colortbl}

% ---- Colours ----
\usepackage{xcolor}
\definecolor{notebg}{HTML}{FFF3CD}
\definecolor{noteborder}{HTML}{FFC107}
\definecolor{boxred}{HTML}{F8D7DA}
\definecolor{boxredborder}{HTML}{F5C6CB}

% ---- Boxes ----
\usepackage{mdframed}
\usepackage{tcolorbox}
\tcbuselibrary{most}

\newtcolorbox{keybox}[1][]{
  colback=blue!5,
  colframe=blue!70!black,
  arc=4pt,
  boxrule=1pt,
  fonttitle=\bfseries,
  title={#1}
}

\newtcolorbox{warnbox}[1][]{
  colback=boxred,
  colframe=boxredborder,
  coltitle=black!75,
  arc=4pt,
  boxrule=1pt,
  fonttitle=\bfseries,
  title={#1}
}

\newtcolorbox{notebox}[1][]{
  colback=notebg,
  colframe=noteborder,
  coltitle=black!75,
  arc=4pt,
  boxrule=1pt,
  fonttitle=\bfseries,
  title={#1}
}

% ---- Hyperlinks & TOC ----
\usepackage[hidelinks]{hyperref}
\usepackage{bookmark}

% ---- Enumerate ----
\usepackage{enumitem}

% ---- Miscellaneous ----
\usepackage{icomma}
\usepackage{siunitx}
\usepackage{textcomp}
\usepackage[bottom]{footmisc}

% ---- Diagramas (gráficas y esquemas) ----
\usepackage{pgfplots}
\pgfplotsset{compat=1.18}
\usetikzlibrary{babel}  % babel español activa `>`; esto evita romper las flechas `->`
% courses/ElecMag2024/Clases/assets/tikzstyles.tex
% ---------------------------------------------------------------------------
% Estilos compartidos para los diagramas del curso Electromagnetismo (2024).
% Fuente única del "look" visual (colores, grosores, escalas) para que todas
% las figuras (TikZ / pgfplots / CircuiTikz) sean consistentes entre clases.
%
% Es una COPIA de courses/Fisica3-2015/Clases/assets/tikzstyles.tex: mismo
% dominio, misma paleta. Copia y no symlink para que cada curso sea
% autocontenido y la paleta de Física III pueda evolucionar aparte.
%
% Uso: la clase debe cargar antes `tikz` y `pgfplots` (y `circuitikz` si la
%      clase tiene circuitos), y luego % courses/ElecMag2024/Clases/assets/tikzstyles.tex
% ---------------------------------------------------------------------------
% Estilos compartidos para los diagramas del curso Electromagnetismo (2024).
% Fuente única del "look" visual (colores, grosores, escalas) para que todas
% las figuras (TikZ / pgfplots / CircuiTikz) sean consistentes entre clases.
%
% Es una COPIA de courses/Fisica3-2015/Clases/assets/tikzstyles.tex: mismo
% dominio, misma paleta. Copia y no symlink para que cada curso sea
% autocontenido y la paleta de Física III pueda evolucionar aparte.
%
% Uso: la clase debe cargar antes `tikz` y `pgfplots` (y `circuitikz` si la
%      clase tiene circuitos), y luego % courses/ElecMag2024/Clases/assets/tikzstyles.tex
% ---------------------------------------------------------------------------
% Estilos compartidos para los diagramas del curso Electromagnetismo (2024).
% Fuente única del "look" visual (colores, grosores, escalas) para que todas
% las figuras (TikZ / pgfplots / CircuiTikz) sean consistentes entre clases.
%
% Es una COPIA de courses/Fisica3-2015/Clases/assets/tikzstyles.tex: mismo
% dominio, misma paleta. Copia y no symlink para que cada curso sea
% autocontenido y la paleta de Física III pueda evolucionar aparte.
%
% Uso: la clase debe cargar antes `tikz` y `pgfplots` (y `circuitikz` si la
%      clase tiene circuitos), y luego \input{../assets/tikzstyles.tex}
% (compilar desde el directorio de la clase para que la ruta relativa resuelva).
% ---------------------------------------------------------------------------

% ---- Paleta (alineada con las cajas del preámbulo: keybox/notebox/warnbox) ----
\definecolor{figblue}{HTML}{1F4E79}   % azul principal (líneas, keybox)
\definecolor{figred}{HTML}{C0392B}    % rojo de énfasis (warnbox)
\definecolor{figamber}{HTML}{B8860B}  % ámbar (notebox)
\definecolor{figgray}{HTML}{555555}   % auxiliares, ejes, guías

% ---- CircuiTikz: escala y grosor de los componentes ----
% Guarda \ifdefined: la electrostática (Clases 1-14) no carga `circuitikz`, y
% sin la guarda el \input revienta con `Undefined control sequence \ctikzset`.
% Cuando el curso llegue a circuitos, el bloque se activa solo.
\ifdefined\ctikzset
  \ctikzset{
    bipoles/thickness=1,
    resistors/scale=0.9,
    capacitors/scale=0.9,
  }
\fi

% ---- TikZ: estilos reutilizables ----
% Nota: las puntas se declaran como `-{Latex[...]}` y no como `->`. La forma
% con llaves NO contiene un `>` literal, así que es inmune al `>` activo que
% introduce babel español (ver CLAUDE.md §7.3.3).
\usetikzlibrary{arrows.meta,calc,angles,quotes}

\tikzset{
  figline/.style={figblue, line width=0.9pt},
  figaux/.style={figgray, dashed, line width=0.6pt},
  % vectores
  figvec/.style={figblue, line width=0.9pt, -{Latex[length=2.2mm,width=1.6mm]}},
  figvecres/.style={figred, line width=1.3pt, -{Latex[length=2.6mm,width=1.9mm]}},
  figvecaux/.style={figgray, line width=0.7pt, -{Latex[length=1.8mm,width=1.3mm]}},
  figaxis/.style={figgray, line width=0.6pt, -{Latex[length=2mm,width=1.4mm]}},
  % cargas puntuales
  figcarga/.style={circle, inner sep=0pt, minimum size=5.4pt, draw=none},
  figpos/.style={figcarga, fill=figred},
  figneg/.style={figcarga, fill=figblue},
  figneutra/.style={figcarga, fill=figgray},
  % etiquetas
  figlbl/.style={font=\small},
  figlblaux/.style={font=\footnotesize, figgray},
}

% ---- pgfplots: estilo base de las gráficas ----
\pgfplotsset{
  emfig/.style={
    width=11cm, height=6.5cm,
    axis lines=left,
    xlabel style={font=\small},
    ylabel style={font=\small},
    tick label style={font=\footnotesize},
    every axis plot/.append style={line width=1pt},
    grid=major,
    grid style={figgray!20},
    legend cell align=left,
    legend style={font=\footnotesize, draw=figgray!40},
  },
}

% (compilar desde el directorio de la clase para que la ruta relativa resuelva).
% ---------------------------------------------------------------------------

% ---- Paleta (alineada con las cajas del preámbulo: keybox/notebox/warnbox) ----
\definecolor{figblue}{HTML}{1F4E79}   % azul principal (líneas, keybox)
\definecolor{figred}{HTML}{C0392B}    % rojo de énfasis (warnbox)
\definecolor{figamber}{HTML}{B8860B}  % ámbar (notebox)
\definecolor{figgray}{HTML}{555555}   % auxiliares, ejes, guías

% ---- CircuiTikz: escala y grosor de los componentes ----
% Guarda \ifdefined: la electrostática (Clases 1-14) no carga `circuitikz`, y
% sin la guarda el \input revienta con `Undefined control sequence \ctikzset`.
% Cuando el curso llegue a circuitos, el bloque se activa solo.
\ifdefined\ctikzset
  \ctikzset{
    bipoles/thickness=1,
    resistors/scale=0.9,
    capacitors/scale=0.9,
  }
\fi

% ---- TikZ: estilos reutilizables ----
% Nota: las puntas se declaran como `-{Latex[...]}` y no como `->`. La forma
% con llaves NO contiene un `>` literal, así que es inmune al `>` activo que
% introduce babel español (ver CLAUDE.md §7.3.3).
\usetikzlibrary{arrows.meta,calc,angles,quotes}

\tikzset{
  figline/.style={figblue, line width=0.9pt},
  figaux/.style={figgray, dashed, line width=0.6pt},
  % vectores
  figvec/.style={figblue, line width=0.9pt, -{Latex[length=2.2mm,width=1.6mm]}},
  figvecres/.style={figred, line width=1.3pt, -{Latex[length=2.6mm,width=1.9mm]}},
  figvecaux/.style={figgray, line width=0.7pt, -{Latex[length=1.8mm,width=1.3mm]}},
  figaxis/.style={figgray, line width=0.6pt, -{Latex[length=2mm,width=1.4mm]}},
  % cargas puntuales
  figcarga/.style={circle, inner sep=0pt, minimum size=5.4pt, draw=none},
  figpos/.style={figcarga, fill=figred},
  figneg/.style={figcarga, fill=figblue},
  figneutra/.style={figcarga, fill=figgray},
  % etiquetas
  figlbl/.style={font=\small},
  figlblaux/.style={font=\footnotesize, figgray},
}

% ---- pgfplots: estilo base de las gráficas ----
\pgfplotsset{
  emfig/.style={
    width=11cm, height=6.5cm,
    axis lines=left,
    xlabel style={font=\small},
    ylabel style={font=\small},
    tick label style={font=\footnotesize},
    every axis plot/.append style={line width=1pt},
    grid=major,
    grid style={figgray!20},
    legend cell align=left,
    legend style={font=\footnotesize, draw=figgray!40},
  },
}

% (compilar desde el directorio de la clase para que la ruta relativa resuelva).
% ---------------------------------------------------------------------------

% ---- Paleta (alineada con las cajas del preámbulo: keybox/notebox/warnbox) ----
\definecolor{figblue}{HTML}{1F4E79}   % azul principal (líneas, keybox)
\definecolor{figred}{HTML}{C0392B}    % rojo de énfasis (warnbox)
\definecolor{figamber}{HTML}{B8860B}  % ámbar (notebox)
\definecolor{figgray}{HTML}{555555}   % auxiliares, ejes, guías

% ---- CircuiTikz: escala y grosor de los componentes ----
% Guarda \ifdefined: la electrostática (Clases 1-14) no carga `circuitikz`, y
% sin la guarda el \input revienta con `Undefined control sequence \ctikzset`.
% Cuando el curso llegue a circuitos, el bloque se activa solo.
\ifdefined\ctikzset
  \ctikzset{
    bipoles/thickness=1,
    resistors/scale=0.9,
    capacitors/scale=0.9,
  }
\fi

% ---- TikZ: estilos reutilizables ----
% Nota: las puntas se declaran como `-{Latex[...]}` y no como `->`. La forma
% con llaves NO contiene un `>` literal, así que es inmune al `>` activo que
% introduce babel español (ver CLAUDE.md §7.3.3).
\usetikzlibrary{arrows.meta,calc,angles,quotes}

\tikzset{
  figline/.style={figblue, line width=0.9pt},
  figaux/.style={figgray, dashed, line width=0.6pt},
  % vectores
  figvec/.style={figblue, line width=0.9pt, -{Latex[length=2.2mm,width=1.6mm]}},
  figvecres/.style={figred, line width=1.3pt, -{Latex[length=2.6mm,width=1.9mm]}},
  figvecaux/.style={figgray, line width=0.7pt, -{Latex[length=1.8mm,width=1.3mm]}},
  figaxis/.style={figgray, line width=0.6pt, -{Latex[length=2mm,width=1.4mm]}},
  % cargas puntuales
  figcarga/.style={circle, inner sep=0pt, minimum size=5.4pt, draw=none},
  figpos/.style={figcarga, fill=figred},
  figneg/.style={figcarga, fill=figblue},
  figneutra/.style={figcarga, fill=figgray},
  % etiquetas
  figlbl/.style={font=\small},
  figlblaux/.style={font=\footnotesize, figgray},
}

% ---- pgfplots: estilo base de las gráficas ----
\pgfplotsset{
  emfig/.style={
    width=11cm, height=6.5cm,
    axis lines=left,
    xlabel style={font=\small},
    ylabel style={font=\small},
    tick label style={font=\footnotesize},
    every axis plot/.append style={line width=1pt},
    grid=major,
    grid style={figgray!20},
    legend cell align=left,
    legend style={font=\footnotesize, draw=figgray!40},
  },
}


% ---- Title ----
\title{\textbf{Clase 19: Potencial vector magnético} \\[0.3em]
        \large{Desarrollo multipolar, momento dipolar magnético y campo de un dipolo}}
\author{Transcripto y expandido --- Curso Electromagnetismo (OpenFING)}
\date{}

\begin{document}

\maketitle
\thispagestyle{empty}
\newpage

\tableofcontents
\newpage

\section{Un ejemplo más de Ampère: cable grueso}

Repaso: la clase pasada se probó que, para corrientes estacionarias
($\div \vec J = 0$), Biot y Savart implica Ampère; eso cierra un asunto
que en Física III «había que creer o reventar».

Ejemplo de recordatorio. Un cable muy largo (se omiten los efectos de borde en
las puntas) de radio $R$, perpendicular al pizarrón, con densidad de corriente
$J$ \textbf{uniforme}.

Mismas simetrías que el cable infinito: invariancia por rotaciones alrededor
del eje, de modo que las líneas de campo ---cerradas e invariantes--- son
circunferencias, y el módulo de $\vec B$ es constante sobre cada una. Tomando
una de radio $r$ como curva amperiana, esté adentro o afuera del cable:

\[
B\,2\pi r = \mu_0\,I_{\text{adentro}},
\qquad
I_{\text{adentro}} =
\begin{cases}
J\,\pi r^2 & r<R\\[2pt]
J\,\pi R^2 & r>R
\end{cases}
\]

\begin{keybox}
\[
\boxed{\ B(r) =
\begin{cases}
\dfrac{\mu_0 J}{2}\,r & r<R\\[8pt]
\dfrac{\mu_0 J}{2}\,\dfrac{R^2}{r} & r>R
\end{cases}\ }
\]
\end{keybox}

\begin{figure}[H]
  \centering
  % Clase 19 — módulo del campo de un cable grueso de radio R con densidad de
% corriente uniforme, en unidades de B_max = mu_0 J R / 2:
% lineal adentro (r < R), como 1/r afuera (r > R).
\begin{tikzpicture}
  \begin{axis}[
    emfig, width=11cm, height=6cm,
    xmin=0, xmax=4.2, ymin=0, ymax=1.2,
    xtick={0,1,2,3,4}, ytick={0,0.5,1},
    xlabel={$r/R$}, ylabel={$B/B_{\max}$},
  ]
    \addplot[figaux, forget plot] coordinates {(1,0) (1,1.2)};
    \addplot[figblue, domain=0:1, samples=2, forget plot] {x};
    \addplot[figblue, domain=1:4.2, samples=80, forget plot] {1/x};
    \node[font=\small, anchor=west] at (axis cs:0.08,0.55) {$\propto r$};
    \node[font=\small, anchor=west] at (axis cs:2.2,0.62) {$\propto 1/r$};
    \node[font=\footnotesize, text=figgray, anchor=west] at (axis cs:1.05,1.1) {borde del cable};
  \end{axis}
\end{tikzpicture}

  \caption{Módulo del campo de un cable de radio $R$ con densidad de corriente uniforme, en unidades de su valor máximo, $\mu_0 J R/2$.}
\end{figure}

El módulo crece linealmente hasta $r=R$, donde alcanza su máximo
$B_{\max} = \mu_0 J R/2$, y afuera decrece como $1/r$.

Con esto el campo magnético tiene, como el eléctrico, sus dos leyes:

\begin{center}
\renewcommand{\arraystretch}{1.35}
\small
\begin{tabularx}{\textwidth}{l>{\raggedright\arraybackslash}X>{\raggedright\arraybackslash}X}
\toprule
 & \textbf{Divergencia} & \textbf{Rotacional} \\
\midrule
$\vec B$ & $\div \vec B = 0$ (no hay monopolos) & $\curl \vec B = \mu_0\vec J$ (Ampère) \\
\bottomrule
\end{tabularx}
\end{center}

\section{Potencial vector magnético}

\subsection{Por qué no hay potencial escalar}

En electrostática el campo se escribe como menos el gradiente de un potencial.
Para $\vec B$ en general no se puede: como $\curl \vec B = \mu_0\vec J$
no es cero, $\vec B$ \textbf{no deriva de un potencial escalar}.

\begin{notebox}
Sí se puede definir un potencial escalar magnético en regiones \textbf{sin
corriente}. Es una situación excepcional, pero aparece en el práctico y más
adelante en el teórico.
\end{notebox}

Hay una propiedad alternativa: la divergencia de $\vec B$ sí es cero. Eso
permite asociarle un potencial que, en lugar de un escalar, es un \textbf{vector}.

\subsection{Biot y Savart como un rotacional}

Se parte de

\[
\vec B(\vec r_2) = \frac{\mu_0}{4\pi}\int_V
\vec J(\vec r_1)\times\frac{\vec r_2-\vec r_1}{|\vec r_2-\vec r_1|^3}\,dV_1
\]
y se busca escribirlo no como un gradiente sino como un rotacional.

\textbf{Paso 1.} El cociente es, salvo constantes, el campo de una carga puntual en
$\vec r_1$, que es menos el gradiente de su potencial:

\[
\vnabla_2\frac{1}{|\vec r_2-\vec r_1|} = -\frac{\vec r_2-\vec r_1}{|\vec r_2-\vec r_1|^3}
\quad\Longrightarrow\quad
\vec B(\vec r_2) = -\frac{\mu_0}{4\pi}\int_V \vec J(\vec r_1)\times
\vnabla_2\frac{1}{|\vec r_2-\vec r_1|}\,dV_1 .
\]

\textbf{Paso 2.} Rotacional de un escalar por un vector (otra vez «la derivada de un
producto»; lo dudoso son los signos, que se verifican a mano o con
Levi-Civita):

\[
\curl (f\,\vec v) = f\,\curl \vec v + \vnabla f\times\vec v .
\]

Con $f = 1/|\vec r_2-\vec r_1|$ y $\vec v = \vec J(\vec r_1)$, el primer
término es nulo porque $\vec J(\vec r_1)$ \textbf{no depende de $\vec r_2$}:

\[
\vnabla_2\times\frac{\vec J(\vec r_1)}{|\vec r_2-\vec r_1|}
= \vnabla_2\frac{1}{|\vec r_2-\vec r_1|}\times\vec J(\vec r_1).
\]

\textbf{Paso 3.} Eso es el integrando de arriba con los factores en el otro orden;
el producto vectorial es anticonmutativo y el signo absorbe el menos. Y el
rotacional, que es una derivada respecto de la variable exterior, \textbf{sale de la
integral}:

\[
\vec B(\vec r_2) = \vnabla_2\times\left[\frac{\mu_0}{4\pi}\int_V
\frac{\vec J(\vec r_1)}{|\vec r_2-\vec r_1|}\,dV_1\right].
\]

\subsection{Definición y analogía con el potencial escalar}

\begin{keybox}
\[
\boxed{\ \vec A(\vec r_2) = \frac{\mu_0}{4\pi}\int_V
\frac{\vec J(\vec r_1)}{|\vec r_2-\vec r_1|}\,dV_1,
\qquad \vec B = \curl \vec A\ }
\]
\end{keybox}

$\vec A$ es el \textbf{potencial vector magnético}. La expresión es casi idéntica a
la del potencial de Coulomb:

\[
\phi(\vec r_2) = \frac{1}{4\pi\varepsilon_0}\int_V
\frac{\rho(\vec r_1)}{|\vec r_2-\vec r_1|}\,dV_1 .
\]

\begin{center}
\renewcommand{\arraystretch}{1.35}
\small
\begin{tabularx}{\textwidth}{l>{\raggedright\arraybackslash}X}
\toprule
\textbf{Eléctrico} & \textbf{Magnético} \\
\midrule
$1/\varepsilon_0$ & $\mu_0$ \\
$\rho$ & $\vec J$ \\
$\phi$ (escalar) & $\vec A$ (vector) \\
$\vec E = -\vnabla\phi$ & $\vec B = \curl \vec A$ \\
\bottomrule
\end{tabularx}
\end{center}

\textbf{¿Qué se ganó?} La objeción es razonable: se pasó de un vector a otro vector,
no de tres componentes a una como en el caso eléctrico. Igual se gana bastante:

\begin{itemize}
  \item la integral ---que es lo que da trabajo--- es \textbf{mucho más simple}: no tiene
  productos vectoriales, y el rotacional se toma una sola vez al final;
  \item es mnemotécnicamente cómoda: en Biot y Savart uno duda del orden de los
  factores y de los signos; ésta se recuerda a partir de la de Coulomb.
\end{itemize}

\subsection{Qué es exacto y qué es aproximado}

La deducción se apoya en Biot y Savart, que es una ley \textbf{aproximada}. Pero:

\begin{notebox}[Resultado matemático (enunciado, no demostrado)]
Todo campo vectorial que
verifique $\div \vec B = 0$ se puede escribir como
$\vec B = \curl \vec A$ para algún campo vectorial $\vec A$. Es análogo
a que un campo de rotacional nulo se escribe como un gradiente, con hipótesis
del mismo tipo sobre la región (simplemente conexa, etc.); se toma todo el
espacio «para no hacernos lío».
\end{notebox}

Como hasta donde se sabe los experimentos nunca detectaron monopolos,
$\div \vec B = 0$ es exacta, y entonces:

\begin{center}
\renewcommand{\arraystretch}{1.35}
\small
\begin{tabularx}{\textwidth}{>{\raggedright\arraybackslash}X>{\raggedright\arraybackslash}X}
\toprule
\textbf{Afirmación} & \textbf{Estatus} \\
\midrule
$\vec B = \curl \vec A$ para algún $\vec A$ & \textbf{Exacta}; no depende de Biot y Savart \\
La expresión integral concreta de $\vec A$ & Basada en Biot y Savart; tendrá correcciones cuando se corrija esa ley \\
\bottomrule
\end{tabularx}
\end{center}

Resumen del docente: (1) es útil, simplifica cuentas; (2) no tanto como en el
caso eléctrico, porque siguen siendo tres componentes; (3) $\vec B$ siempre es
el rotacional de un potencial vector, pero la fórmula disponible está al nivel
de aproximación de Biot y Savart; (4) se parece mucho a la del potencial
escalar.

\section{Desarrollo multipolar magnético}

Pregunta: un conjunto de corrientes localizadas en una región; se quiere el
campo magnético \textbf{lejos} de ella. Igual que en el caso eléctrico, se hace un
desarrollo en el tamaño de la región, y conviene empezar por el potencial.
Cada componente de $\vec A$ obedece la misma expresión que $\phi$, así que la
cuenta es la misma, con particularidades.

\subsection{Planteo para un circuito}

Se trata el caso más común, corriente $I$ en un circuito $C$. Con la regla
$\vec J\,dV\to I\,d\vec l$:

\[
\vec A(\vec r_2) = \frac{\mu_0 I}{4\pi}\oint_C
\frac{d\vec l_1}{|\vec r_2-\vec r_1|} .
\]

Se toma el origen cerca del circuito, de modo que $|\vec r_1|$ sea del orden
del tamaño $a$ del circuito, y se estudia $|\vec r_2|\gg|\vec r_1|$.

\begin{figure}[H]
  \centering
  % Clase 19 — planteo del desarrollo multipolar magnético. El origen se toma
% cerca del circuito, de modo que |r_1| es del orden de su tamaño a, y se mira
% el campo en un punto lejano, |r_2| >> a.
\begin{tikzpicture}[scale=1]

  % circuito
  \draw[figblue, line width=1.3pt]
    plot[smooth cycle, tension=0.8] coordinates {(-1.0,-0.1) (-0.5,0.85) (0.6,1.0) (1.2,0.2) (0.7,-0.8) (-0.4,-0.85)};
  \draw[figvecres] (-0.05,-0.97) -- (0.5,-0.9);
  \node[figlbl, text=figred] at (0.15,-1.3) {$I$};
  \node[figlbl] at (-1.25,0.75) {$C$};

  % origen y r_1
  \fill[figgray] (0,0) circle (1.4pt);
  \node[figlblaux] at (-0.22,-0.22) {$O$};
  \draw[figvec] (0,0) -- (0.5,0.97);
  \node[figlbl] at (-0.08,0.6) {$\vec r_1$};

  % tamaño a
  \draw[figgray, line width=0.6pt, {Latex[length=1.6mm]}-{Latex[length=1.6mm]}] (-1.05,-1.75) -- (1.25,-1.75);
  \node[figlblaux] at (0.1,-2.05) {$\sim a$};

  % punto lejano
  \coordinate (P) at (8.2,2.1);
  \draw[figvec] (0,0) -- (P);
  \node[figlbl] at (4.4,0.75) {$\vec r_2$};
  \draw[figaux, -{Latex[length=2mm,width=1.4mm]}] (0.5,0.97) -- (P);
  \node[figlbl] at (3.6,1.9) {$\vec r_2-\vec r_1$};
  \fill[figblue] (P) circle (1.6pt);
  \node[figlbl, anchor=west] at (7.0,2.6) {$|\vec r_2| \gg a$};

\end{tikzpicture}

  \caption{El origen se toma cerca del circuito, de modo que $|\vec r_1|$ es del orden de su tamaño $a$, y se mira el campo en un punto lejano.}
\end{figure}

\subsection{Desarrollo de \texorpdfstring{$1/|\vec r_2-\vec r_1|$}{1/|r2 - r1|}}

El mismo cálculo del caso eléctrico:

\[
\frac{1}{|\vec r_2-\vec r_1|}
= \frac{1}{\left(r_2^2 + r_1^2 - 2\,\vec r_1\cdot\vec r_2\right)^{1/2}}
= \frac{1}{r_2}\left(1 - \frac{2\,\vec r_1\cdot\vec r_2}{r_2^2}
+ \frac{r_1^2}{r_2^2}\right)^{-1/2}.
\]

\begin{center}
\renewcommand{\arraystretch}{1.35}
\small
\begin{tabularx}{\textwidth}{l>{\raggedright\arraybackslash}X>{\raggedright\arraybackslash}X}
\toprule
\textbf{Aproximación} & \textbf{Orden en $r_1/r_2$} & \textbf{Significado} \\
\midrule
Monopolar & 0 & Se reemplaza $\vec r_1$ por 0: todo el circuito en el origen \\
Dipolar & 1 & --- \\
Cuadrupolar & 2 & --- \\
\bottomrule
\end{tabularx}
\end{center}

En el caso eléctrico se llegó al cuadrupolo; acá alcanza con el \textbf{dipolar}.
Con $(1+x)^\alpha = 1+\alpha x+\mathcal O(x^2)$ y $\alpha = -\tfrac12$ (el
término $r_1^2/r_2^2$ ya es de segundo orden):

\[
\frac{1}{|\vec r_2-\vec r_1|} = \frac{1}{r_2}\left[1 +
\frac{\vec r_1\cdot\vec r_2}{r_2^2} + \mathcal O\!\left(\frac{r_1^2}{r_2^2}\right)\right].
\]

Sustituyendo (con $d\vec l_1 = d\vec r_1$):

\[
\vec A(\vec r_2) = \frac{\mu_0 I}{4\pi}\left[
\frac{1}{r_2}\oint_C d\vec r_1
+ \frac{1}{r_2^3}\oint_C d\vec r_1\,(\vec r_1\cdot\vec r_2)
+ \mathcal O\!\left(\frac{1}{r_2^3}\right)\right].
\]

El primer término decrece como $1/r_2$, el segundo como $1/r_2^2$ (hay un
$r_2$ arriba) y el cuadrupolar, despreciado, como $1/r_2^3$.

\subsection{El término monopolar es cero}

\[
\oint_C d\vec r_1 = 0 :
\]
la integral de línea sobre una curva cerrada se anula. \textbf{No hay término
monopolar}, como tenía que ser: no existe el equivalente magnético de la
carga. El campo magnético de un circuito decrece, a lo sumo, dipolarmente.

\subsection{El término dipolar}

Hay que calcular $\oint_C d\vec r_1\,(\vec r_1\cdot\vec r_2)$, y se quiere
sacar $\vec r_2$ de la integral, donde está entreverado en un producto escalar.

\begin{notebox}
Este tramo del video es una \textbf{regrabación}: el docente avisa que la
explicación original «fue un enredo inútil» y la rehace prolija.
\end{notebox}

\textbf{Ecuación (1): triple producto vectorial.}

\[
(\vec r_1\times d\vec r_1)\times\vec r_2
= -\,\vec r_1\,(d\vec r_1\cdot\vec r_2) + d\vec r_1\,(\vec r_1\cdot\vec r_2).
\]

\textbf{Ecuación (2): un diferencial exacto.} Variando $\vec r_1$ con $\vec r_2$
fijo:

\[
d\!\left[\vec r_1\,(\vec r_2\cdot\vec r_1)\right]
= d\vec r_1\,(\vec r_1\cdot\vec r_2) + \vec r_1\,(\vec r_2\cdot d\vec r_1).
\]

La expresión buscada aparece en las dos, y la otra aparece con signos opuestos.
\textbf{Sumando}:

\[
(\vec r_1\times d\vec r_1)\times\vec r_2 + d\!\left[\vec r_1\,(\vec r_2\cdot\vec r_1)\right]
= 2\,d\vec r_1\,(\vec r_1\cdot\vec r_2).
\]

Integrando sobre la curva cerrada, el diferencial exacto da cero (se vuelve al
punto de partida) y $\vec r_2$ sale de la integral:

\[
\oint_C d\vec r_1\,(\vec r_1\cdot\vec r_2)
= \left[\frac12\oint_C \vec r_1\times d\vec r_1\right]\times\vec r_2 .
\]

\subsection{Momento dipolar magnético y potencial del dipolo}

El corchete es el \textbf{vector área} del circuito, el mismo que apareció al
calcular el torque sobre una espira (Clase 16):

\[
\vec{\mathcal A} = \frac12\oint_C \vec r_1\times d\vec r_1 .
\]

(Con $\mathcal A$ cursiva para no confundirlo con el potencial vector.)
Multiplicado por la corriente es el \textbf{momento dipolar magnético}:

\[
\vec m = I\,\vec{\mathcal A} = \frac{I}{2}\oint_C \vec r_1\times d\vec r_1 .
\]

Para una curva plana, $\vec{\mathcal A}$ es el de Física III: módulo igual al
área, dirección perpendicular al plano, sentido según la mano derecha. Para un
circuito en el espacio es más complicado, pero la integral lo da.

\begin{keybox}
\[
\boxed{\ \vec A(\vec r) = \frac{\mu_0}{4\pi}\,\frac{\vec m\times\vec r}{r^3}
+ \mathcal O\!\left(\frac{1}{r^3}\right)\ }
\]
\end{keybox}

Es la \textbf{aproximación dipolar magnética} del potencial vector. Comparada con el
dipolo eléctrico:

\[
\phi(\vec r) = \frac{1}{4\pi\varepsilon_0}\,\frac{\vec p\cdot\vec r}{r^3} + \dots
\]
el producto escalar pasa a ser vectorial, $1/\varepsilon_0$ pasa a $\mu_0$ y
$\vec p$ pasa a $\vec m$.

\section{Campo magnético de un dipolo}

\subsection{Cálculo del rotacional}

$\vec B = \curl \vec A$, con $\vec m$ \textbf{constante} (depende del dipolo,
no del punto de observación). Usando el rotacional de un producto vectorial
(Clase 18), sólo sobreviven los términos donde las derivadas caen sobre
$\vec r/r^3$:

\[
\curl \left(\vec m\times\frac{\vec r}{r^3}\right)
= \vec m\;\div \!\left(\frac{\vec r}{r^3}\right)
- (\vec m\cdot\vnabla)\frac{\vec r}{r^3}.
\]

\textbf{Primer término.} $\vec r/r^3$ es el campo de una carga puntual: su
divergencia es $4\pi\,\delta(\vec r)$, nula fuera del origen.

\begin{warnbox}[La delta en el origen no interesa]
La aproximación dipolar es para $r$
grande; en $r=0$ nada de esto vale y habría que mirar el detalle de la
distribución de corriente. Se descarta el término.
\end{warnbox}

\textbf{Segundo término.} Componente a componente. Derivando arriba y abajo:

\[
m_x\,\pdv{x}\,\frac{\vec r}{r^3}
= \frac{m_x\,\hat\imath}{r^3} - \frac{3\,m_x\,x\;\vec r}{r^5},
\]
y lo mismo con $y$ y $z$. Sumando los tres:

\[
(\vec m\cdot\vnabla)\frac{\vec r}{r^3}
= \frac{\vec m}{r^3} - \frac{3\,(\vec m\cdot\vec r)\,\vec r}{r^5}.
\]

\subsection{Resultado y cuándo falla}

Con el signo menos de adelante (que el docente se había «morfado» y corrige):

\begin{keybox}
\[
\boxed{\ \vec B(\vec r) = \frac{\mu_0}{4\pi}\left[
\frac{3\,(\vec m\cdot\vec r)\,\vec r}{r^5} - \frac{\vec m}{r^3}\right]\ }
\]
\end{keybox}

Es el \textbf{campo magnético dipolar}. Es casi idéntico al campo de un dipolo
eléctrico: se pasa de uno a otro cambiando $1/\varepsilon_0\to\mu_0$ y
$\vec p\to\vec m$. El nombre «dipolar» es totalmente apropiado.

\begin{notebox}
\textbf{Cualquier circuito, visto de lejos, se comporta así}, siempre que
$\vec m\neq0$. En el caso eléctrico, para llegar al cuadrupolo tenían que
anularse carga y momento dipolar; acá el término de carga se anula solo.
\end{notebox}

¿Cuándo $\vec m = 0$? Una sola curva plana nunca lo tiene nulo. Sí puede
anularse, por ejemplo, con dos circuitos girando en sentidos opuestos, o con
una curva que sube y vuelve por el otro lado: es como sumar dos dipolos
opuestos.

\section{Flujo magnético por superficies abiertas}

Definición que hará falta más adelante, para la ley de Faraday. El flujo de
$\vec B$ por una superficie \textbf{cerrada} es cero; pero interesa también el flujo
por superficies abiertas.

Dada una curva cerrada orientada $C$ y dos superficies $S_1$ y $S_2$ con ese
mismo borde, cada una orientada por la \textbf{regla de la mano derecha} (la
convención del curso):

\[
\Phi_{S_1} = \int_{S_1}\vec B\cdot\hat n\,dS,
\qquad
\Phi_{S_2} = \int_{S_2}\vec B\cdot\hat n\,dS .
\]

\begin{figure}[H]
  \centering
  % Clase 19 — dos superficies S_1 y S_2 con el mismo borde C. Las dos se
% orientan con la regla de la mano derecha respecto del sentido de C, así que
% sus normales apuntan "hacia el mismo lado". Juntas forman una superficie
% cerrada; con normal saliente hay que invertir la de S_2.
\begin{tikzpicture}[scale=1]

  % S_2 (por abajo)
  \fill[figgray!14] (-2.6,0) .. controls (-2.3,-2.3) and (2.3,-2.3) .. (2.6,0) arc (0:-180:2.6 and 0.8);
  \draw[figaux] (-2.6,0) .. controls (-2.3,-2.3) and (2.3,-2.3) .. (2.6,0);

  % S_1 (por arriba)
  \fill[figblue!12] (-2.6,0) arc (180:360:2.6 and 0.8) .. controls (2.3,2.3) and (-2.3,2.3) .. (-2.6,0);
  \draw[figline] (2.6,0) .. controls (2.3,2.3) and (-2.3,2.3) .. (-2.6,0);

  % curva C
  \draw[figaux] (2.6,0) arc (0:180:2.6 and 0.8);
  \draw[figblue, line width=1.3pt] (-2.6,0) arc (180:360:2.6 and 0.8);
  \draw[figvecres] (-0.3,-0.8) -- (0.4,-0.8);
  \node[figlbl] at (2.95,-0.35) {$C$};

  % normales
  \draw[figvecaux] (1.3,1.28) -- (1.6,2.0);
  \node[figlblaux] at (1.9,2.05) {$\hat n$};
  \draw[figvecaux] (1.3,-1.42) -- (1.5,-0.72);
  \node[figlblaux] at (1.8,-0.82) {$\hat n$};

  \node[figlbl] at (-2.05,1.4) {$S_1$};
  \node[figlbl] at (-2.05,-1.5) {$S_2$};

  % lineas de campo
  \foreach \bx in {-1.1,-0.4} {
    \draw[figvec] (\bx,-2.6) -- (\bx,2.7);
  }
  \node[figlbl] at (-0.05,2.6) {$\vec B$};

\end{tikzpicture}

  \caption{Dos superficies con el mismo borde $C$, orientadas las dos con la regla de la mano derecha. Las líneas de campo que cruzan una cruzan la otra.}
\end{figure}

Son iguales. Intuitivamente: el flujo cuenta líneas de campo, y todas las que
pasan por una pasan por la otra. Formalmente: la unión $S_1\cup S_2$ es una
superficie cerrada, pero al tomar la normal saliente hay que invertir la de una
de las dos:

\[
0 = \Phi_{S_1\cup S_2} = \Phi_{S_1} - \Phi_{S_2}
\quad\Longrightarrow\quad
\boxed{\ \Phi_{S_1} = \Phi_{S_2}\ }
\]

\textbf{El flujo magnético depende sólo del borde}, no de la superficie elegida para
calcularlo.

\begin{notebox}
Con el campo eléctrico no pasa lo mismo: si entre las dos superficies queda
una carga, por Gauss el flujo por la superficie cerrada no es cero. En el
caso magnético, como no hay monopolos, vale siempre.
\end{notebox}

\section{Adelanto: Laplace fuera del electromagnetismo}

Quedan ocho minutos y el docente los usa para anunciar algo de la clase
siguiente. La ecuación de Laplace ya apareció dos veces (electrostática y
corrientes estacionarias en medios óhmicos), y aparecerá una tercera en
problemas magnéticos: donde $\vec J = 0$ es $\curl \vec B = 0$, se escribe
$\vec B$ como un gradiente y, con $\div \vec B = 0$, vuelve Laplace.

Pero es mucho más general: vale siempre que haya un fenómeno \textbf{lineal,
homogéneo e isotrópico, sin fuentes}, determinado por lo que pasa en el borde.
Dos ejemplos que se anuncian y \textbf{no se desarrollan} en esta clase:

\begin{center}
\renewcommand{\arraystretch}{1.35}
\small
\begin{tabularx}{\textwidth}{>{\raggedright\arraybackslash}X>{\raggedright\arraybackslash}X>{\raggedright\arraybackslash}X}
\toprule
\textbf{Ejemplo} & \textbf{Qué obedece Laplace} & \textbf{Dimensión} \\
\midrule
Membrana tensa de un tambor (de borde cualquiera) deformada \textbf{un poquito} por un taco, en equilibrio & La altura de la membrana como función de $(x,y)$ & Plano \\
Salón cerrado con una estufa, en régimen estacionario & El campo de temperaturas & Espacio \\
\bottomrule
\end{tabularx}
\end{center}

Para deformaciones grandes de la membrana el problema es no lineal y más
complicado; fuera del régimen estacionario, la temperatura obedece otra
ecuación, la del calor.

\noindent\emph{La próxima clase dedica un rato a esas aplicaciones de la ecuación de Laplace
en otras ramas de la física.}

\end{document}
